% !TEX root = ../../proposal.tex

\section{Phase 1: Characterising the Composability Gap}
\label{sec:phase_1}

As discussed in Section~\ref{sec:compositional_diffusion}, Bayes-style score composition does not define a truly out-of-distribution joint distribution. Its success therefore depends on the model already containing a usable approximation to the desired joint configuration within its learned training distribution; when that assumption fails, score-level composition can deviate substantially from the behaviour of a model conditioned on a single semantic description of the same scene.

This leads to a central problem:
\begin{center}
{\setlength{\fboxsep}{6pt}%
\fbox{%
\parbox{\dimexpr 0.96\linewidth - 2\fboxsep - 2\fboxrule\relax}{%
\textbf{Logical composition and semantic composition are not equivalent, and their discrepancy depends on the structure of the underlying data distribution.}%
}%
}}
\end{center}

% In particular, joint compositions form complex, highly structured distributions that may be difficult to represent or sample from, especially when independence or factorization assumptions are violated \citep{du2023reduce,du2024compositional}.

In this phase, we empirically characterise this discrepancy as a \textit{composability gap}, and study how it varies across different types of concept pairs. To do so, we construct controlled groups of compositions designed to span a range of regimes from approximately factorized pairs to strongly conflicting or entangled concepts. This grouping allows us to probe the structural conditions under which logical composition succeeds or fails, and to identify the mechanisms underlying this gap.

This phase tests Hypothesis~1 in three steps:
\begin{itemize}
    \item We establish that the composability gap exists and varies systematically across concept pairs chosen to span a range of conditions, from manifold-supported co-occurrence to adversarial concept collision.
    \item We characterise the gap dynamically by tracking full denoising trajectories, allowing us to identify when the two generation processes diverge and how that divergence develops over time.
    \item We test whether the endpoint reached by logical composition can be recovered through language conditioning alone. Persistent failure here supports the claim that logical composition reaches targets that standard semantic prompting does not.
\end{itemize}

\subsection{Concept Pair Taxonomy}
\label{subsec:taxonomy}

The taxonomy is grounded in the conditions under which \citet{mechanisms2025projective} predicts composition to succeed or fail. Rather than treating successes and failures as isolated examples, Phase~1 groups concept pairs into interpretable regimes so that H1 can be tested against a clear structural ordering.

We partition concept pairs into four groups:
\begin{enumerate}

    \item \textbf{Co-occurring Concepts (Group~1):}
    These pairs frequently co-occur in the training distribution, so the model is likely to have already learned a usable joint representation. This serves as a positive control: logical and semantic composition should remain closely aligned, yielding the smallest gaps.

    \item \textbf{Feature-Space Disentanglement (Group~2):}
    These pairs act on largely distinct feature dimensions, such as object--style combinations. They form a second positive control: even without strong co-occurrence in the training data, composition should remain reliable when the factors are structurally compatible.

    \item \textbf{Feature Overlap (Group~3):}
    These pairs have weak joint support and no clear feature-space separation. They are expected to produce intermediate behaviour: larger and more variable gaps than the positive controls, but not necessarily universal failure.


    \item \textbf{Adversarial Collision (Group~4):}
    These pairs compete for the same semantic slot, object role, or spatial allocation. They represent the strongest violation of the conditions required for successful composition, and should therefore produce the largest and most systematic failures.
    % If the ordering $\text{G1} \approx \text{G2} < \text{G3} < \text{G4}$ is observed, this supports the claim that the composability gap is governed by the geometry of concept representations, rather than by a uniform limitation of the composition operator.

\end{enumerate}
This taxonomy defines the pair regimes used in the analyses below.

% Figures~\ref{fig:multimodel_grid_g1}--\ref{fig:multimodel_grid_g4} show decoded outputs from three representative pairs per group, compared across three model scales: SD~1.4, SDXL, and SD~3.5 Medium.
% Within each figure, rows are organised by model (three rows per concept pair) and columns show the four conditions: solo $c_1$, solo $c_2$, semantic composition, and logical composition (logical composition/AND).
% Because all conditions within a row share the same $x_T$, differences between the semantic and logical columns are attributable to the composition operator alone.
% The multi-model format rules out model-specific implementation details as a confound and demonstrates that failure patterns are a structural property of the concept pairs rather than an artefact of any single architecture.

% \begin{figure}[H]
%     \centering
%     \includegraphics[width=\textwidth]{chapters/research_method/media/group1_cooccurrence_multimodel_grid.png}
%     \caption{\textbf{Group~1 (Co-occurring Concepts)} across SD~1.4, SDXL, and SD~3.5~Medium. Each block of three rows shows one concept pair (camel/desert landscape, butterfly/flower meadow, lion/savanna at sunset); columns show solo $c_1$, solo $c_2$, semantic composition, and logical composition (PoE). Both operators converge to visually comparable outputs across all models: the semantic and logical columns produce similarly coherent scenes, and neither hybridisation nor concept dominance is observed. Consistency across all three architectures confirms that Group~1 serves as a reliable positive control.}
%     \label{fig:multimodel_grid_g1}
% \end{figure}

% \begin{figure}[H]
%     \centering
%     \includegraphics[width=\textwidth]{chapters/research_method/media/group2_disentangled_multimodel_grid.png}
%     \caption{\textbf{Group~2 (Feature-Space Disentangled)} across three model scales. Pairs combine a concrete object with a rendering style (dog/oil-painting, lighthouse/watercolour, bicycle/sketch). Logical composition closely tracks the semantic column across all models, consistent with the feature-space disentanglement prediction of \citet{bradley2025} (Theorem~6.1): when composed concepts act along approximately orthogonal feature dimensions, PoE can be satisfied without off-manifold compromise. Replication across architectures confirms that the disentanglement is a property of the concept pair, not of any single model.}
%     \label{fig:multimodel_grid_g2}
% \end{figure}

% \begin{figure}[H]
%     \centering
%     \includegraphics[width=\textwidth]{chapters/research_method/media/group3_ood_multimodel_grid.png}
%     \caption{\textbf{Group~3 (Feature Overlap, OOD)} across three model scales. Pairs have intermediate CLIP-proxy alignment (bathtub/streetlamp, black grand piano/white vase, typewriter/cactus). Logical composition produces inconsistent results: some pairs and models yield coherent co-presence, while others exhibit partial concept dominance or incoherent blending. The behaviour is pair-dependent and model-dependent within the group, consistent with the prediction that moderate orthogonality violations lead to degraded but non-catastrophic composition. The semantic column remains coherent for all pairs and models.}
%     \label{fig:multimodel_grid_g3}
% \end{figure}

% \begin{figure}[H]
%     \centering
%     \includegraphics[width=\textwidth]{chapters/research_method/media/group4_collision_multimodel_grid.png}
%     \caption{\textbf{Group~4 (Adversarial Collision)} across three model scales. Pairs share a single semantic slot and have high CLIP-proxy alignment (cat/dog, cat/owl, cat/bear). Logical composition (PoE) systematically produces either \emph{hybridisation}---a chimeric entity blending features of both concepts---or \emph{concept dominance}---one concept suppresses the other. The semantic column renders both concepts as distinct, co-present agents across all models. Critically, both failure modes appear in all three architectures: the pattern is not an artefact of any single model's scale or implementation, but a structural consequence of the concepts competing over the same latent degrees of freedom, providing strong preliminary evidence for H1.}
%     \label{fig:multimodel_grid_g4}
% \end{figure}

% Two failure modes are consistently observed in Groups~3 and~4.
% The first is \textbf{hybridisation}: logical composition collapses two concepts into a single incoherent entity rather than negotiating a shared scene.
% In cat/owl pairs (Group~4), PoE/AND frequently produces a single animal blending feline and avian features, whereas the monolithic prompt renders a cat and an owl as distinct co-present agents.
% This is consistent with both concepts competing for the same spatial coordinates in the absence of a joint training attractor.
% The second failure mode is \textbf{concept dominance}: one concept's score field systematically overwhelms the other, reducing the suppressed concept to a distorted artefact or removing it entirely.
% Crucially, both failure modes are group-predictable under the Factorised Conditionals framework: Groups~1 and~2 rarely exhibit either, Group~4 almost always exhibits at least one, and Group~3 exhibits intermediate behaviour consistent with the ``harder but not categorically failed'' prediction.


\subsection{Latent Trajectory Dynamics}
\label{subsec:trajectories}
The purpose of analysing the latent trajectory dynamics is to determine whether the gap between logical and semantic composition appears only at the decoded endpoint or is already present during denoising. 
To do this, we compare semantic composition (e.g generating complex through a single prompt) and logical composition under a shared-noise protocol.
The analysis will be organised by taxonomy group so that H1 is evaluated with respect to structured compositional regimes rather than a pooled average over heterogeneous pairs.

For each taxonomy group, we select six concept pairs and run twenty-four shared-noise seeds per pair. Each pair-seed record uses the same initial latent state for solo $c_1$, solo $c_2$, the monolithic prompt, and PoE. Seeds are treated as repeated measurements within a pair, so pair-level summaries remain the unit of group comparison.

This analysis asks three questions: whether the two trajectories remain close or diverge, when that divergence appears during denoising, and whether the expected Group~1--4 ordering is visible in the terminal gap. If the largest and earliest divergences concentrate in the structurally hardest groups, then H1 is supported at the level of latent dynamics rather than only decoded outputs.

This sets up the reachability analysis below.


\subsection{Reachability Analysis}
\label{subsec:reachability}

In this phase, we ask whether the terminal state of logical composition can be reproduced by prompt-conditioned generation. For each pair-seed record, the final composed image is inverted into a continuous prompt text embedding using SD-IPC~\citep{clip_inverter}, and the model is rerun from the same initial noise under that recovered conditioning.

If the regenerated trajectory returns to logical composition's region, the gap may reflect a limitation of prompt choice within the semantic manifold. If it instead drifts back toward nearby semantic or dominant-concept modes, that supports the stronger claim that logical composition reaches targets that standard prompting does not readily recover. 
We compare this behaviour across the same Group~1--4 taxonomy, using trajectory closeness, terminal distance to the original PoE run, and output-level similarity as the main indicators.

\subsection{Summary}
\label{subsec:p1_predictions}

Phase~1 characterises the composability gap as a structured empirical phenomenon. The taxonomy
identifies the concept regimes in which the gap should be minimal or severe, the trajectory analysis
tests whether semantic and logical composition separate during denoising, and the reachability
analysis tests whether the endpoint reached by logical composition remains recoverable through
prompt-conditioned generation. Taken together, these analyses support H1 when the gap worsens
predictably from Groups~1--2 to Group~4, with the largest and earliest divergence appearing in the
hardest regimes.

What Phase~1 does not resolve is the cause of this pattern. Because Stable Diffusion \citep{rombach2022high} fixes both the internal
representations and the score compositional method, these results cannot yet determine whether the gap
is fundamentally a representation problem or a compositional method. Phase~2 addresses that question
directly by testing the representation hypothesis in a controlled setting, asking whether our novel architecture that
progressively disentangles latents reduces the same compositional failures observed here.
